\documentclass[11pt,a4paper]{article}
\usepackage[utf8]{inputenc}
\usepackage[T1]{fontenc}
\usepackage[french]{babel}
\usepackage{lmodern}
\usepackage[margin=2.3cm]{geometry}
\usepackage{amsmath,amssymb}
\usepackage{booktabs}
\usepackage{tabularx}
\usepackage{enumitem}
\usepackage{xcolor}
\usepackage[hidelinks]{hyperref}

\newcommand{\code}[1]{\texttt{#1}}

\title{ManageVet\\[0.3em]\large Simulation économique de clinique vétérinaire\\Document de conception --- version 0.1}
\author{}
\date{Octobre 2026}

\begin{document}
\maketitle
\tableofcontents

\section{Concept}

Chaque joueur (étudiant vétérinaire) dirige une clinique pendant une année
simulée découpée en quatre trimestres. Un trimestre de jeu correspond à une
semaine réelle : les décisions sont saisies au fil de la semaine et stockées
\emph{en attente} ; une tâche planifiée unique résout ensuite le trimestre
pour toutes les cliniques d'un même serveur. Aucun calcul en temps réel n'est
nécessaire et la micro-gestion est évitée.

Une dizaine de cliniques partagent une carte commune et se disputent une
patientèle simulée, qui arbitre selon le prix, la distance, l'offre de
services et la réputation.

\section{Architecture du dépôt}

\begin{tabularx}{\textwidth}{@{}lX@{}}
\toprule
Application & Rôle \\
\midrule
\code{config} & Configuration Django (paramètres lus dans l'environnement). \\
\code{game} & Serveur de jeu, carte et zones de population, catalogue de référence, vues, commandes de gestion. \\
\code{clinics} & Clinique, prestations et prix pratiqués, personnel, matériel (achat ou leasing). \\
\code{inventory} & Gammes de produits, stocks, réassort automatique, valorisation au CMUP. \\
\code{accounting} & Extrait du PCG, écritures en partie double, rapports trimestriels. \\
\code{decisions} & Décisions en attente et leur application à la résolution. \\
\code{engine} & Moteur : modèle de demande (pur Python) et résolution d'un trimestre. \\
\bottomrule
\end{tabularx}

\medskip
Le module \code{engine.demand} ne dépend pas de Django : il est testable
isolément et réutilisable pour des simulations hors ligne.

\section{Modèle de données}

\begin{description}[style=nextline]
\item[\code{GameServer}] Partie : dimensions de la grille, trimestre courant,
  nombre total de trimestres, échéance de la prochaine résolution, graine
  aléatoire, surcharges de paramètres (JSON) pour les stress tests.
\item[\code{PopulationZone}] Une case de la grille : nombre de chiens et de
  chats, sensibilité au prix.
\item[\code{Clinic}] Position, propriétaire (un joueur par serveur),
  trésorerie, loyer, honoraires de l'expert-comptable, réputation
  $r \in [0,1]$.
\item[\code{ServiceType} / \code{ClinicService}] Catalogue des actes (prix de
  référence, fréquence par animal, temps vétérinaire et ASV, consommables,
  matériel requis) et prix pratiqué par chaque clinique.
\item[\code{StaffMember}] Vétérinaire ou ASV, salaire brut mensuel,
  trimestres d'entrée et de sortie.
\item[\code{EquipmentType} / \code{ClinicEquipment}] Matériel acquis par achat
  (immobilisé et amorti) ou par leasing (charge trimestrielle).
\item[\code{ProductLine} / \code{StockItem}] Gamme de produits ; stock d'une
  clinique avec quantité, CMUP, niveau de réassort cible et prix de vente.
\item[\code{PendingDecision}] Type, charge utile JSON, statut (en attente,
  appliquée, rejetée, annulée), message.
\item[\code{JournalEntry} / \code{JournalLine}] Écriture en partie double ;
  l'équilibre débit/crédit est vérifié à la saisie.
\item[\code{QuarterReport}] Synthèse figée d'un trimestre : compte de
  résultat recalculé depuis le journal, indicateurs d'activité.
\end{description}

\paragraph{Invariant.} La trésorerie \code{Clinic.cash} est une valeur
dénormalisée, modifiée uniquement par la passation d'écritures ; elle est en
permanence égale au solde du compte 512. Ce point est vérifié par les tests.

\section{Carte et positionnement}

La carte est une grille abstraite de $W \times H$ cases (par défaut
$20 \times 20$, une case $= 2$~km). La population d'une case $(x,y)$ est
la somme d'un fond rural et de quelques villes gaussiennes de centres
$(c_x^k, c_y^k)$ et de taille $s_k$ :
\[
  N(x,y) = N_{\text{rural}}\,u + N_{\text{ville}} \sum_k s_k
  \exp\!\left(-\frac{(x-c_x^k)^2 + (y-c_y^k)^2}{2\sigma^2}\right),
  \qquad u \sim \mathcal{U}(0{,}6\,;\,1{,}4).
\]
La proportion de chats augmente en ville ; la sensibilité au prix est plus
forte à la campagne. Chaque clinique occupe une case distincte ; les
distances sont euclidiennes. Ce choix suffit à créer des marchés urbains
denses et concurrentiels et des zones rurales peu servies, sans données
géographiques réelles.

\section{Moteur de résolution}

\subsection{Ordre des étapes}

\begin{enumerate}
  \item Application des décisions en attente (chacune dans un point de
    sauvegarde : une décision rejetée ne laisse aucune trace).
  \item Réassort automatique jusqu'au niveau cible (compte 607).
  \item Répartition de la patientèle entre toutes les cliniques.
  \item Pour chaque clinique : actes plafonnés par la capacité (706 et 6022),
    ventes plafonnées par le stock (707), salaires et charges sociales (641,
    645), loyer (6132), honoraires (6226), autres charges (628), leasing
    (6122), dotations (6811), variation de stock (6037), agios (6615).
  \item Mise à jour de la réputation et rapport trimestriel.
  \item Passage au trimestre suivant ou fin de partie.
\end{enumerate}

L'ensemble est exécuté dans une seule transaction. Le hasard provient d'un
générateur initialisé par (graine, trimestre, position de la clinique) : la
résolution est reproductible.

\subsection{Choix de la clinique : modèle logit}

Pour une zone $z$ et une clinique $c$ :
\[
  U_{zc} = \alpha - \beta_d\, d_{zc} - \beta_p\,(I_c - 1)\,\sigma_z
  + \beta_o\,(O_c - 1) + \beta_r\,(r_c - 0{,}5)
\]
où $d_{zc}$ est la distance en km, $I_c$ l'indice de prix de la clinique,
$\sigma_z$ la sensibilité au prix de la zone, $O_c \in [0,1]$ la part
pondérée des services proposés et $r_c$ la réputation. Avec une option
extérieure d'utilité $U_0$ (ne pas consulter ce trimestre) :
\[
  s_{zc} = \frac{e^{U_{zc}}}{e^{U_0} + \sum_{c'} e^{U_{zc'}}},
  \qquad
  A_c = \mu \sum_z N_z\, s_{zc},
\]
$\mu$ étant un multiplicateur global de demande. Les poids de l'indice de
prix et du ratio d'offre sont ceux du panier de référence,
$w_k = p_k^{\text{réf}} f_k$ :
\[
  I_c = \frac{\sum_{k \in \mathcal{S}_c} w_k\, p_{ck}/p_k^{\text{réf}}}{\sum_{k \in \mathcal{S}_c} w_k},
  \qquad
  O_c = \frac{\sum_{k \in \mathcal{S}_c} w_k}{\sum_k w_k}.
\]

\subsection{Actes et capacité}

Le volume demandé pour l'acte $k$ vaut
$q_{ck} = A_c\, f_k\, (p_{ck}/p_k^{\text{réf}})^{-\varepsilon_s}\,(1+\eta)$,
avec $\eta$ un aléa uniforme de faible amplitude. La capacité est le temps
disponible du personnel ; le taux de service est
\[
  \kappa_c = \min\!\left(1,\;
  \frac{n^{\text{véto}}_c\, H^{\text{véto}}}{\sum_k q_{ck}\, t^{\text{véto}}_k},\;
  \frac{n^{\text{ASV}}_c\, H^{\text{ASV}}}{\sum_k q_{ck}\, t^{\text{ASV}}_k}\right)
\]
et les actes réalisés sont $\kappa_c\, q_{ck}$. La demande non servie est
perdue en V1 (elle n'est pas redirigée vers les concurrents).

\subsection{Ventes de produits et stocks}

Les ventes de la gamme $g$ sont proportionnelles aux animaux effectivement
reçus et plafonnées par le stock disponible :
\[
  v_{cg} = \min\!\left(Q_{cg},\;
  \kappa_c\, A_c^{(\text{espèce})}\, u_g\,(p_{cg}/p_g^{\text{conseillé}})^{-\varepsilon_p}\right).
\]
Le réassort recomplète le stock jusqu'au niveau cible en début de
trimestre ; le stock est valorisé au CMUP. La variation de stock
$\Delta S = S_{\text{initial}} - S_{\text{final}}$ est passée au compte 6037,
de sorte que le coût des marchandises vendues vaut $607 + 6037$.

\subsection{Réputation}

\[
  r_c^{t+1} = \lambda\, r_c^t + (1-\lambda)\,\bigl(0{,}7\,\kappa_c + 0{,}3\,\phi_c\bigr),
\]
où $\phi_c$ est le taux de service du stock et $\lambda$ l'inertie.

\section{Formats d'établissement}

Chaque clinique choisit un format, modifiable par une décision du
trimestre. Cette décision est appliquée \emph{après} les autres : les
recrutements et le matériel décidés le même trimestre comptent pour les
conditions d'accès.

\begin{center}
\begin{tabularx}{\textwidth}{@{}lXX@{}}
\toprule
Format & Conditions et coûts & Effets \\
\midrule
Standard & --- & Offre selon le matériel et les spécialistes présents. \\
Low-cost & Offre limitée aux actes sans matériel ni spécialiste
  (consultation, vaccination, stérilisation, détartrage). &
  Temps par acte $\times 0{,}75$, consommables $\times 0{,}85$ ;
  ratio d'offre $O_c$ plus faible. \\
CHV & Voir ci-dessous. &
  Bonus d'utilité $b_c = 0{,}6$ ; urgences et hospitalisation continue,
  recrutées sur un bassin régional. \\
\bottomrule
\end{tabularx}
\end{center}

\subsection{Centre hospitalier vétérinaire}

Les conditions reprennent l'arrêté du 13 mars 2015 relatif aux catégories
d'établissements de soins vétérinaires et le cahier des charges « CHV pour
animaux de compagnie » de l'Ordre des vétérinaires (24 mars 2016), tels que
résumés par les sources citées en fin de document :

\begin{itemize}
  \item au moins 6 vétérinaires et 6 personnes qualifiées à temps plein ;
    seules les ASV titrées (échelon 5) comptent comme personnes qualifiées ;
  \item un vétérinaire et une personne qualifiée présents sur place
    24\,h/24, 7\,j/7, et ouverture au public en permanence ;
  \item au moins 6 vétérinaires spécialistes au sens de l'article R.\,242-34
    du code rural et de la pêche maritime, tous à l'échelon 5 (cadre
    spécialisé), dans au moins 2 disciplines parmi médecine interne,
    chirurgie et imagerie. Leur activité est restreinte aux actes de leur
    discipline, mais ils reçoivent les cas référés que les autres cliniques
    ne peuvent pas prendre en charge ;
  \item au moins trois techniques d'imagerie parmi radiographie,
    échographie, scanner et IRM ;
  \item un laboratoire : analyseurs de biochimie, d'hématologie et
    d'ionogrammes.
\end{itemize}

Le premier passage en CHV impose des travaux d'aménagement de
60\,000 euros (compte 2135, amortis sur 40 trimestres, payés une seule
fois). Les conditions sont revérifiées à chaque trimestre ; une clinique
qui ne les remplit plus redevient standard (événement signalé dans son
rapport).

\paragraph{Coût de la permanence.} Hors des heures d'ouverture de jour
(8\,h--20\,h du lundi au samedi), il reste $96$\,h par semaine, soit
$H_p = 1\,248$\,h par trimestre, à couvrir par un vétérinaire et une
personne qualifiée. Avec $1\,607$\,h annuelles par ETP, cela mobilise
$H_p / 401{,}75 \approx 3{,}1$ ETP par poste. L'activité de nuit étant
faiblement facturable (productivité $\rho = 0{,}3$), les heures retirées,
$3{,}1\,(1-\rho) \times 220$\,h par trimestre, sont prises d'abord sur la
réserve des généralistes, puis, à défaut, sur celle des spécialistes
(au prorata).
Les heures de nuit (22\,h--7\,h), de dimanche et de jours fériés,
$H_m \approx 1\,060$\,h par trimestre et par poste, sont majorées selon la
convention collective :
\[
  M = H_m \left( \bar{s}_{\text{véto}}\, m_{\text{véto}} + \bar{s}_{\text{ASV}}\, m_{\text{ASV}} \right),
  \qquad \bar{s} = \frac{\text{salaire brut mensuel moyen}}{151{,}67},
\]
avec $m_{\text{véto}} = 20\,\%$ et $m_{\text{ASV}} = 15\,\%$ (compte 6411,
charges sociales en sus). S'y ajoutent 8\,000 euros par trimestre de
fonctionnement 24\,h/24 (énergie, maintenance, sécurité ; compte 628).

\section{Personnel et convention collective}

Les rémunérations suivent la convention collective des établissements de
soins vétérinaires : IDCC 1875 pour le personnel non vétérinaire,
annexe VII (ex-IDCC 2564) pour les vétérinaires salariés. Le salaire
minimum mensuel brut d'un échelon vaut
\[
  s_{\min} = \text{coefficient} \times \text{valeur du point},
  \qquad \text{point 2026} = 17{,}96~\text{euros (avenant n° 91)}.
\]

\begin{center}
\begin{tabular}{@{}llrr@{}}
\toprule
Poste & Échelon & Coefficient & Minimum 2026 \\
\midrule
Vétérinaire & 2 -- cadre débutant & 152 & 2\,729,92 \\
Vétérinaire & 3 -- cadre confirmé A & 182 & 3\,268,72 \\
Vétérinaire & 4 -- cadre confirmé B & 210 & 3\,771,60 \\
Vétérinaire & 5 -- cadre spécialisé & 240 & 4\,310,40 \\
Auxiliaire & 3 -- auxiliaire vétérinaire qualifié & 110 & 1\,975,60 \\
Auxiliaire & 4 -- auxiliaire vétérinaire confirmé & 113 & 2\,029,48 \\
Auxiliaire & 5 -- auxiliaire spécialisé vétérinaire & 120 & 2\,155,20 \\
\bottomrule
\end{tabular}
\end{center}

Un vétérinaire spécialiste relève obligatoirement de l'échelon 5 ; le
joueur fixe librement un salaire au moins égal au minimum (le marché
rémunère souvent les spécialistes bien au-delà).

\subsection{Effets de l'échelon}

La convention fixe des minima de rémunération, pas des niveaux de
productivité : les effets suivants sont des hypothèses de jeu,
paramétrables.

\begin{enumerate}
  \item \textbf{Productivité.} Les heures facturables d'un salarié sont
    multipliées par un coefficient $p_e$ selon son échelon :
    \begin{center}
    \begin{tabular}{@{}lcccc@{}}
    \toprule
    Vétérinaire, échelon & 2 & 3 & 4 & 5 \\
    $p_e$ & 0,80 & 1,00 & 1,10 & 1,15 \\
    \midrule
    Auxiliaire, échelon & & 3 & 4 & 5 \\
    $p_e$ & & 0,90 & 0,95 & 1,00 \\
    \bottomrule
    \end{tabular}
    \end{center}
    Un cadre débutant coûte 17\,\% de moins qu'un confirmé A au minimum
    conventionnel, mais produit 20\,\% de moins.
  \item \textbf{Qualité perçue.} L'expérience moyenne de l'équipe,
    $\bar{x} \in [-1, 1]$ (chaque salarié est situé entre le plus bas et le
    plus haut échelon de son poste), s'ajoute à la satisfaction :
    \[
      \text{satisfaction} = 0{,}7\,\kappa_c + 0{,}3\,\phi_c + 0{,}15\,\bar{x},
    \]
    bornée à $[0, 1]$, ce qui fait évoluer la réputation.
  \item \textbf{Qualification.} Seules les ASV titrées (échelon 5) sont des
    personnes qualifiées au sens des conditions du CHV.
\end{enumerate}

\subsection{Réserves de temps et spécialistes}

Le temps vétérinaire est réparti en réserves : une réserve
\emph{généraliste} (vétérinaires sans spécialité) et une réserve par
spécialité. Chaque acte consomme le temps de sa réserve (les actes de
spécialité ne peuvent être réalisés que par un spécialiste de la
discipline, et un spécialiste ne réalise que ceux-là) et le temps ASV,
commun à tous les actes. La part servie d'un acte est
$\min(\kappa_{\text{réserve}}, \kappa_{\text{ASV}})$. Les actes de
spécialité exigent aussi leur matériel : bloc de chirurgie spécialisée
(instrumentation orthopédique, amplificateur de brillance), colonne
d'endoscopie, scanner ou échographe pour l'imagerie. Ils sont interdits en
low-cost. La valeur du point et les taux de majoration sont des paramètres
de simulation (revalorisation, passage à la convention fusionnée IDCC 3256
annoncée pour 2027).

\section{Difficulté de recrutement}

Une décision de recrutement est une \emph{tentative}. En début de
résolution, toutes les tentatives du serveur sont tranchées ensemble.

\paragraph{Chances individuelles.} Pour une offre de salaire $s$ d'une
clinique $c$ sur un profil $k$ (généraliste, spécialiste, auxiliaire) :
\[
  p = \operatorname{borne}_{[0{,}02\,;\,0{,}98]}\!\left(
    b_k \cdot \mu \cdot \min\!\left(\left(\frac{s}{s^{\text{marché}}_k}\right)^{\eta}, \bar{f}\right)
    \cdot \left(1 - \lambda\,(1-\delta_c)\right)
    \cdot \left(1 + \omega\,(2 r_c - 1)\right)
    \cdot \gamma_{k,c}
  \right)
\]
\begin{itemize}
  \item $b_k$ : chance de base au salaire du marché (généraliste 0,60,
    spécialiste 0,35, auxiliaire 0,85) ; $\mu$ : multiplicateur de pénurie
    (1 par défaut, 0,5 pour une forte pénurie) ;
  \item $s^{\text{marché}}_k = s_{\min} \times (1 + \pi_k)$, avec des
    primes de marché de 15\,\% (généralistes), 45\,\% (spécialistes) et
    5\,\% (auxiliaires) au-dessus du minimum conventionnel ; élasticité
    $\eta = 3$, plafond $\bar{f} = 1{,}6$ ;
  \item $\delta_c \in [0, 1]$ : densité de la patientèle autour de la
    clinique rapportée à la case la plus dense ; pénalité rurale
    $\lambda = 0{,}35$ ;
  \item $r_c$ : réputation, poids $\omega = 0{,}2$ ;
  \item $\gamma_{k,c}$ : effet du format CHV, $+20\,\%$ pour les spécialistes
    (plateau technique), $-10\,\%$ pour les généralistes (gardes de nuit).
\end{itemize}

\paragraph{Vivier et concurrence.} Chaque trimestre, le marché du serveur
compte un nombre limité de candidats : 5 vétérinaires généralistes,
10 auxiliaires et 1 spécialiste par discipline. Parmi les offres
acceptées, les plus attractives ($p$ multiplié par un aléa de $\pm 20\,\%$)
l'emportent dans la limite du vivier ; les autres échouent. Constituer
l'équipe de 6 spécialistes d'un CHV prend donc plusieurs trimestres, et
les cliniques se disputent les mêmes profils.

\paragraph{Frais.} Chaque tentative coûte des frais d'annonce ou de
cabinet (1\,500 euros pour un généraliste, 6\,000 pour un spécialiste,
500 pour un auxiliaire ; compte 6231), dus même en cas d'échec. Une offre
invalide (sous le minimum conventionnel) est rejetée sans frais.

L'écran de recrutement affiche le salaire du marché et les chances
estimées en direct. Le salaire proposé par défaut est celui du marché.

\section{Urgences et cas référés}

Les urgences, l'hospitalisation continue, les actes de spécialistes et le
scanner sont recrutés sur un \emph{bassin régional} : les propriétaires,
ou leur vétérinaire traitant, acceptent d'aller plus loin. Pour chaque acte
de ce type, la demande de toutes les zones est répartie entre les seules
cliniques qui le proposent, par le même logit que la patientèle, avec une
sensibilité à la distance réduite ($\beta^{\text{réf}}_d = 0{,}08$ par km
au lieu de $0{,}30$), sans effet du prix ni de l'offre généraliste. Ces
actes n'entrent pas dans le panier de référence de $I_c$ et $O_c$.

\paragraph{Calibrage.} Avec les valeurs par défaut, le CHV de la démo
(6 spécialistes, 8 généralistes, 8 ASV titrées, seul CHV du territoire)
réalise environ 550\,000 à 660\,000 euros de chiffre d'affaires par
trimestre ; sa marge nette passe de 6\,\% à 19\,\% au fil de l'année, à
mesure que sa réputation progresse, pour une masse salariale de 58 à
70\,\% du chiffre d'affaires. Une clinique standard employant un seul
chirurgien spécialiste et bien placée atteint 31 à 35\,\% de marge : un
chirurgien référent concentre des actes à forte valeur. Ces ordres de
grandeur restent à confronter à des données réelles de la profession.

\section{Emprunts bancaires}

Un emprunt de capital $P$ sur $n$ trimestres au taux annuel $t$ est
remboursé par échéances constantes, à partir du trimestre suivant le
déblocage :
\[
  a = P\,\frac{r}{1-(1+r)^{-n}}, \qquad r = \frac{t}{4}.
\]
Chaque échéance se décompose en intérêts $r \cdot \mathrm{CRD}$ (compte
6611) et amortissement du capital (compte 164). La banque fixe le taux à
la souscription :
\[
  t = t_0 + \pi_\ell\,\frac{D + P}{\mathit{CP}} + \pi_p\,\mathbb{1}_{\text{dernier résultat} < 0},
\]
avec $t_0 = 4{,}5\,\%$, $\pi_\ell = 1$ point par unité d'endettement,
$\pi_p = 1$ point, $D$ la dette existante et $\mathit{CP}$ les capitaux
propres (capital et résultats cumulés). Elle refuse si $\mathit{CP} \le 0$
ou si $D + P > \min(3\,\mathit{CP},\ 400\,000)$. Le remboursement anticipé
solde le capital restant dû et coûte une indemnité de 3\,\% (compte 668).

\section{Comptabilité}

Écritures types (montants HT) :

\begin{center}
\begin{tabular}{@{}lll@{}}
\toprule
Opération & Débit & Crédit \\
\midrule
Apport en capital & 512 & 101 \\
Achat de matériel & 2154 & 512 \\
Dotation aux amortissements (linéaire) & 6811 & 28154 \\
Loyer de leasing & 6122 & 512 \\
Réassort & 607 & 512 \\
Prestations & 512 & 706 \\
Déblocage d'un emprunt & 512 & 164 \\
Échéance d'emprunt & 6611 / 164 & 512 \\
Travaux d'aménagement CHV & 2135 & 512 \\
Majorations de nuit (CHV) & 6411 & 512 \\
Frais de recrutement & 6231 & 512 \\
Ventes de marchandises & 512 & 707 \\
Salaires / charges patronales & 6411 / 645 & 512 \\
Loyer / honoraires / divers & 6132 / 6226 / 628 & 512 \\
Hausse de stock (baisse : inverse) & 37 & 6037 \\
Agios & 6615 & 512 \\
\bottomrule
\end{tabular}
\end{center}

\paragraph{Simplifications assumées en V1.} Pas de TVA (neutre pour une
clinique assujettie) ; règlements au comptant (pas de comptes 401, 411, 421,
431) ; un salarié recruté est payé sur le trimestre entier ; pas de cession
de matériel acheté ; pas d'impôt sur les sociétés ni de clôture annuelle.

\section{Paramètres et stress tests}

Tous les coefficients sont regroupés dans \code{engine/parameters.py}
(classe \code{SimulationParameters}). Un serveur peut en surcharger une
partie via son champ JSON \code{parameters}, par exemple :

\begin{center}
\begin{tabular}{@{}ll@{}}
\toprule
Scénario & Surcharge \\
\midrule
Choc de demande & \code{demand\_multiplier = 0.7} \\
Inflation fournisseurs & \code{purchase\_cost\_multiplier = 1.25} \\
Hausse des charges sociales & \code{employer\_contribution\_rate = 0.50} \\
Clientèle plus sensible au prix & \code{beta\_price = 5} \\
\bottomrule
\end{tabular}
\end{center}

\section{Inscription et choix de l'emplacement}

L'étudiant crée un compte, puis rejoint une partie en choisissant une case
libre de la carte. Pour chaque case, l'écran affiche :

\begin{itemize}
  \item la patientèle potentielle dans un rayon de chalandise
    $R = 6$~km : $P(x,y) = \sum_{d((x,y),z) \le R} N_z$ ;
  \item la distance au concurrent le plus proche et le nombre de concurrents
    dans le rayon ;
  \item le loyer trimestriel, croissant avec la densité :
    \[
      L(x,y) = L_0 + \Delta L \,\frac{P(x,y)}{\max_{(x',y')} P(x',y')},
      \qquad L_0 = 3\,000~\text{euros},\ \Delta L = 3\,000~\text{euros}.
    \]
\end{itemize}

Ce loyer variable crée l'arbitrage pédagogique entre centre-ville (clientèle
nombreuse, loyer élevé, concurrence) et campagne (loyer bas, clientèle
dispersée). Chaque partie a un nombre maximal de cliniques, peut fermer les
inscriptions et peut exiger un code d'accès communiqué par l'enseignant.
L'installation est faite sous verrou du serveur de jeu, ce qui empêche deux
étudiants de prendre simultanément la même case ou la dernière place.

\section{Écran de décisions du joueur}

Le joueur saisit ses décisions sur une page unique
(\code{/clinics/<id>/decisions/}), organisée en blocs : tarifs et offre de
soins, produits et réassort, personnel, matériel. Chaque formulaire crée des
\code{PendingDecision} pour le trimestre en cours ; rien n'est appliqué
avant la résolution.

\begin{itemize}
  \item \textbf{Valeurs planifiées.} Les formulaires affichent l'état actuel
    corrigé des décisions en attente. Une nouvelle saisie portant sur la même
    prestation (ou gamme, ou matériel) annule et remplace la précédente ;
    revenir à la valeur actuelle annule simplement la décision.
  \item \textbf{Impact estimé.} Chaque décision en attente est accompagnée de
    son effet sur la trésorerie du trimestre (salaire chargé sur trois mois,
    prix d'achat ou loyer de leasing, indemnité de départ).
  \item \textbf{Validation.} Les erreurs de saisie (prix négatif, salaire
    sous le minimum) sont signalées immédiatement ; les règles dépendant de
    l'état au moment de la résolution restent vérifiées par le moteur.
  \item \textbf{Confidentialité.} Un joueur n'accède qu'à sa clinique ;
    l'équipe pédagogique (\code{is\_staff}) accède à toutes.
  \item \textbf{Concurrence.} La création d'une décision verrouille le
    serveur de jeu : elle ne peut pas viser un trimestre déjà résolu.
\end{itemize}

\section{Planification}

La commande \code{resolve\_quarters} résout tous les serveurs dont
l'échéance est passée ; il suffit de l'appeler par \code{cron} (ou un timer
systemd) une fois par semaine. Une file de tâches (Celery) n'est pas
nécessaire tant que le calcul d'un serveur reste de l'ordre de la seconde.

\section{Pistes pour les versions suivantes}

\begin{itemize}
  \item Rattachement des étudiants à une promotion et tableau de bord
    enseignant.
  \item Progression d'échelon à l'ancienneté (glissement
    vieillesse-technicité), marché de l'emploi des spécialistes.
  \item Redirection de la demande non servie vers les cliniques voisines.
  \item TVA, comptes de tiers, emprunts bancaires, bilan et clôture annuelle.
  \item Événements aléatoires (épidémie, départ d'un salarié) et tableau de
    bord enseignant comparant les stratégies.
  \item Calibrage des paramètres avec des données réelles de la profession.
\end{itemize}

\section{Sources et réserves}

Le texte officiel de l'arrêté, le cahier des charges de l'Ordre et les
avenants salariaux n'ont pas pu être consultés directement lors de la
rédaction ; les valeurs ci-dessus proviennent de sources secondaires
concordantes, qui divergent parfois (notamment sur le nombre de
spécialistes exigés : un ou deux). Elles sont toutes paramétrables et
doivent être vérifiées par l'équipe pédagogique sur Légifrance et auprès de
l'Ordre avant usage.

\begin{itemize}
  \item Arrêté du 13 mars 2015 relatif aux catégories d'établissements de
    soins vétérinaires :
    \url{https://www.legifrance.gouv.fr/loda/id/JORFTEXT000030395433/}
  \item Cahier des charges CHV pour animaux de compagnie (Ordre des
    vétérinaires) :
    \url{https://www.veterinaire.fr/system/files/files/2021-11/cahier-des-charges-ac-chv.pdf}
  \item Synthèse des cahiers des charges (Le Fil vétérinaire) :
    \url{https://www.lefil.vet/_contenus_dyn/articles/106/src/Comparatif_CahierChargesCabinet_Clinique_CVH_AnxCie.pdf}
  \item Grilles 2026 vétérinaires et ASV : \url{https://salerya.fr/conventions/veterinaire/}
  \item Convention des vétérinaires praticiens salariés (annexe VII) :
    \url{https://www.legisocial.fr/conventions-collectives-nationales/2564-veterinaires-praticiens-salaries}
  \item Gardes et astreintes : \url{https://asvinfos.com/1532-gardes-et-astreintes.html}
\end{itemize}

\end{document}
